\section{Aula 5}

Hoje estou um pouco \Frowny, portanto devo ser um pouco mais suscinto.
Na aula, vimos duas  ferramentas bem legais que aparecem em
percolação. As duas abordam probabilidades de eventos crescentes
no hipercubo $\{0,1\}^{[n]}$. Em todos os teoremas seguintes, como
geralmente percolamos arestas,
nosso hipercubo será da forma
$\{0,1\}^{E}$ onde $E$ é o conjunto de arestas de um grafo $G = (V,E)$ finito.
Começamos pela identidade de Russo, para qual precisamos de uma definição.

\begin{definition}
	Dado $A \subset \{0,1\}^E$ um evento crescente e $\omega \in \{0,
	1\}^E$, dizemos
	que uma aresta $e \in E$ é pivotal em $\omega$ em relação a $A$ se vale que
	\[\omega_e^{+} \in A \quad \& \quad \omega_e^- \not \in A\]
	onde $\omega_e^+$ é alterar $\omega$ somente na coordenada $e$ para
	$1$ e $\omega_e^-$
	é definido análogamente.
\end{definition}

Isso é, a aresta $e$ é pivotal em $\omega$ para $A$ se trocar o
estado de $e$ muda se pertencemos a $A$.
Com isso, conseguimos definir um evento
\[
	\partial_e A = \{e \text{ é pivotal em } A\} = \{\omega: e \text{ é
	pivotal em $\omega$ em relação a } A\}.
\]
Dado $\bm{p} \in [0,1]^{E}$, consideramos o modelo $G_{\bm{p}}$ que é percolar
cada aresta $e$ com probabilidade $\bm{p}_e$. Temos a seguinte identidade

\begin{theorem}
	\label{trm:russo}
	Dado $A$ evento crescente, $e \in E$ e pensando em $\Pr_{\bm{p}}(A)$
	como uma função nas entradas $\bm{p}_e$ de $\bm{p}$, temos
	\begin{equation*}
		\frac{d\Pr_{\bm{p}}(A)}{d\bm{p}_e} = \Pr(\partial_e A).
	\end{equation*}
	Em particular, fixando $\bm{p} = (p,p,\dots, p)$, segue que
	\[
		\frac{d\Pr_p(A)}{dp} = \sum_{e \in E} \Pr(\partial_e A) = \Ex
		\#\{\text{arestas pivotais em A}\}.
	\]
\end{theorem}

\begin{proof}
	O Augusto deu uma prova bem bonitinha probabilistica usando o coupling
	com uniformes. Tem uma mais bruta que é uma conta simples no Bela.
\end{proof}

Depois vimos bem legal da desigualdade de Berg-Kestern para eventos crescentes.
É quase que uma recíproca para desigualdade de Harris. Não vou
escrever a prova aqui,
somente as definições necessárias. Uma prova por indução pode ser
achada no Bela também.

\begin{definition}
	Dado um evento $A$ crescente em $\{0,1\}^E$ e $\omega \in A$,
	dizemos que $I \subset E$ é
	uma testemunha para o evento $A$ de $\omega$ se para todo $\omega'$
	tal que $\omega'|_I = \omega|_I$ temos que $\omega' \in A$.
	Considerando todos tais $\omega$, definimos o evento em que $I$ é uma
	testemunha para $A$.
\end{definition}
Essa definição formal é feia, mas a ideia é bem simples. $I$ é uma testemunha
se ter uma configuração específica em $I$, significa pertencer a $A$.
Por exemplo,
se percolamos $G(n,p)$ e estamos interessados no evento de conter
algum triângulo,
se $\omega$ contém um triângulo $T$, então este serve de testemunha para
o evento $A$. Qualquer
outro $\omega'$ que é igual a $\omega$ em $T$,
pertence também a $A$.

Agora podemos ir mais além. Dadas duas famílias crescentes $A$ e $B$,
definimos o evento $A \square B$
sendo os pontos tais que $A$ e $B$ podem ser disjuntamente testemunhadas.
\[A \square B = \{\exists I,J \subset E, I \cap J = \varnothing, I
\text{ é testemunha de } A, J \text{ é testemunha de } B \}\]
Se $A$ é o evento de conter um $K_3$ em $G(n,p)$ e $B$ é o evento de
conter um $C_4$,
$A \square B$ seria o evento de conter um $K_3$ e um $C_4$ aresta-disjuntos.

Conseguimos uma desigualdade, quase recíproca a Harris
\begin{theorem}[Berg-Kestern]
	\label{trm:bk}
	Sejam $A,B \subset \{0,1\}^E$, temos
	\[
		\Pr(A \square B) \leq \Pr(A)\Pr(B).
	\]
\end{theorem}
\begin{remark}
	Essa desigualdade é bem intuitiva, acontecer disjuntamente parece mais difícil
	que acontecer independentemente
\end{remark}

Provamos o teorema somente para $A$ e $B$ crescentes.
Essa prova é bem legal, vamos revelando as arestas sequencialmente e
ficamos atentos
quando que uma aresta está sendo usada de testemunha para $A$ e $B$
ao mesmo tempo.
A prova que o Augusto fez pode ser encontrada no Bela também - mas
ele fez de maneira
muito muito mais legal em sala.
